\section{Results}
\label{sec:results}

\subsection{Benchmark Comparison}
We evaluated our method on three standard benchmarks. The results are summarized
in Table~\ref{tab:results}.

\begin{table}[h]
\centering
\begin{tabular}{lccc}
\hline
\textbf{Method} & \textbf{MNIST} & \textbf{CIFAR-10} & \textbf{ImageNet} \\
\hline
SGD              & 98.2\%  & 91.3\%  & 74.1\% \\
Adam             & 98.5\%  & 92.1\%  & 75.3\% \\
\textbf{Ours}    & \textbf{98.9\%} & \textbf{93.7\%} & \textbf{76.8\%} \\
\hline
\end{tabular}
\caption{Test accuracy comparison across benchmarks}
\label{tab:results}
\end{table}

\subsection{Convergence Rate}
As predicted by Theorem~\ref{thm:main}, our method exhibits linear convergence
with rate $\kappa = L/\mu$. The empirical convergence matches the theoretical
bound closely:
\[
    \frac{\norm{x_k - x^*}}{\norm{x_0 - x^*}} \approx \left(1 - \frac{1}{\kappa}\right)^k
\]

\subsection{Discussion}
The improvement over standard SGD is consistent across all datasets, confirming
the theoretical advantage of the adaptive step size selection (see
Equations~\ref{eq:update}--\ref{eq:stepsize} and Table~\ref{tab:params}).

Our results align with the findings of Brown et al.~\cite{brown2021}, who
observed similar patterns in the convex setting.
